% =============================================================================
% Relatório de Piloto --- LLM Local (Ollama) no Sistema RAG-ANTT
% Ambiente de desenvolvimento como simulacao do perfil de deploy/operacao
% =============================================================================
\documentclass[12pt, a4paper]{article}

% --- Codificacao e idioma ---
\usepackage[utf8]{inputenc}
\usepackage[T1]{fontenc}
\usepackage[brazil]{babel}

% --- Tipografia e layout ---
\usepackage{lmodern}
\usepackage[top=2.5cm, bottom=2.5cm, left=3cm, right=2.5cm]{geometry}
\usepackage{setspace}
\onehalfspacing
\usepackage{indentfirst}
\setlength{\parindent}{1.25cm}
\usepackage{float}

% --- Tabelas ---
\usepackage{booktabs}
\usepackage{array}
\usepackage{longtable}
\usepackage{tabularx}
\usepackage{multirow}

% --- Figuras e cores ---
\usepackage[dvipsnames]{xcolor}
\usepackage{graphicx}
\graphicspath{{figs/}}
\usepackage{tikz}
\usetikzlibrary{arrows.meta,positioning,fit}

% --- Cores e links ---
\usepackage[colorlinks=true, linkcolor=NavyBlue, citecolor=NavyBlue, urlcolor=NavyBlue]{hyperref}

% --- Listas ---
\usepackage{enumitem}

% --- Cabecalho e rodape ---
\usepackage{fancyhdr}
\pagestyle{fancy}
\fancyhf{}
\fancyhead[L]{\small\textit{Piloto LLM Local --- RAG-ANTT (Ollama/CPU)}}
\fancyhead[R]{\small\thepage}
\fancyfoot[C]{\small\textit{Relatório de Avaliação --- DeepFeed}}
\renewcommand{\headrulewidth}{0.4pt}
\renewcommand{\footrulewidth}{0.4pt}

% --- Matematica ---
\usepackage{amsmath}
\usepackage{amssymb}

% --- Citacao em bloco ---
\usepackage{csquotes}

% --- Listagens ---
\usepackage{listings}
\lstset{
  basicstyle=\ttfamily\small,
  frame=single,
  breaklines=true,
  breakatwhitespace=false,
  columns=fullflexible,
  keepspaces=true,
  xleftmargin=0.5cm,
  framexleftmargin=0.5cm,
  backgroundcolor=\color{gray!5},
  rulecolor=\color{gray!50},
  literate={---}{{\textemdash}}1
           {-->}{{$\rightarrow$}}3
           {<--}{{$\leftarrow$}}3,
}

% --- Titulo ---
\title{%
  \vspace{-1cm}
  \textbf{Relatório de Piloto}\\[0.5cm]
  \Large LLM Local via Ollama (CPU) no Sistema RAG-ANTT\\
  Simulação do Ambiente de Deploy e Limitações de Infraestrutura\\[0.3cm]
  \large Agência Nacional de Transportes Terrestres (ANTT)
}

\author{%
  DeepFeed\\
}

\date{Setembro de 2026}

% =============================================================================
\begin{document}

\maketitle
\thispagestyle{empty}

\vspace{1cm}

\begin{center}
  \fbox{\parbox{0.85\textwidth}{
    \centering
    \vspace{0.3cm}
    \textbf{Classificação:} Relatório técnico de piloto --- Projeto RAG-ANTT\\[0.2cm]
    \small O ambiente de desenvolvimento foi configurado para \textbf{simular o perfil}\\
    de deploy/operação permitido pela ANTT (LLM local via Ollama, inferência\\
    em CPU, sem GPU). Este documento registra essa premissa metodológica e\\
    os gargalos de infraestrutura \textbf{pertinentes a esse perfil}, que limitam\\
    o uso de modelos locais mais robustos (acima de $\sim$7B) na operação prevista.
    \vspace{0.3cm}
  }}
\end{center}

\newpage
\tableofcontents
\newpage

% =============================================================================
\section{Introdução e contexto}
% =============================================================================

O Sistema RAG-ANTT consulta documentos regulatórios da Agência Nacional de Transportes Terrestres (ANTT) por meio de recuperação aumentada (RAG) e geração por modelos de linguagem. No âmbito do ofício SUTEC/ANTT e da arquitetura de referência (\texttt{docs/arquitetura\_rag\_api.tex}), a DeepFeed entrega o núcleo de IA, com previsão de LLM local via Ollama para reduzir dependência de APIs externas e manter dados no perímetro institucional.

\subsection{Premissa metodológica: desenvolvimento como simulação do deploy}

O piloto \textbf{não} foi conduzido como um laboratório de hardware genérico. A estação de desenvolvimento foi deliberadamente operada sob o \textbf{mesmo perfil de restrição} esperado no ambiente de deploy/operação da aplicação na Agência (SUTEC/GETIC e cluster Rancher previsto na arquitetura), enquanto não houver acelerador (GPU) e quotas de memória adequadas:

\begin{itemize}
  \item Inferência \textbf{somente em CPU} (sem GPU útil para o runtime do Ollama);
  \item LLM \textbf{local} via Ollama, no perímetro, sem dependência obrigatória de API externa;
  \item Modelos na faixa praticável sob esse teto (tipicamente 3B--7B quantizados), alinhados ao piso técnico documentado na arquitetura (piloto sem GPU: ordem de 7~cores / 22~GB RAM no cluster-alvo).
\end{itemize}

A restrição \textbf{sem GPU} foi confirmada pela GETIC; a arquitetura de referência formaliza a viabilidade do piloto CPU-only e o dimensionamento mínimo. O notebook usado no desenvolvimento não replica, núcleo a núcleo, a VM de produção, mas \textbf{reproduz o perfil operacional decisivo}: CPU-only, memória finita e modelo local Ollama --- exatamente as condições sob as quais a aplicação deverá rodar até novo provisionamento.

O objetivo deste relatório é explicitar quais limitações de \textbf{infraestrutura do perfil de operação simulado} impedem, no momento, o uso de modelos locais significativamente mais capazes (por exemplo, 13B ou superiores com aceleração por GPU). Não se detalham aqui otimizações de software do pipeline RAG nem incidentes locais de estação que não se transferem para esse perfil de deploy.

\subsection{Objetivos do piloto}

\begin{enumerate}[label=\textbf{\arabic*.}]
  \item Simular, em desenvolvimento, o perfil de deploy/operação CPU-only previsto para a ANTT;
  \item Validar a operação do Ollama nesse perfil (LLM local no fluxo Streamlit do RAG);
  \item Confirmar a viabilidade de inferência local com modelos compatíveis com a memória disponível sob CPU-only;
  \item Caracterizar os gargalos de infraestrutura \textbf{pertinentes ao ambiente de operação simulado};
  \item Fundamentar a necessidade de provisionamento adequado (GPU e/ou mais memória) caso a Agência exija qualidade próxima à de LLMs em nuvem \textbf{sem} saída de dados para APIs externas.
\end{enumerate}

\subsection{Escopo}

\begin{table}[H]
\centering
\small
\caption{Escopo do piloto versus arquitetura-alvo}
\label{tab:escopo_piloto}
\begin{tabularx}{\textwidth}{
  >{\raggedright\arraybackslash}X
  >{\raggedright\arraybackslash}p{3.6cm}
  >{\raggedright\arraybackslash}p{4.0cm}
}
\toprule
\textbf{Item} & \textbf{Piloto (simulação)} & \textbf{Arquitetura-alvo} \\
\midrule
Perfil CPU-only / sem GPU & Simulado em desenvolvimento & Previsto até provisionar GPU \\
Provedor Ollama na interface Streamlit & Sim & Sim (cliente da API) \\
Modelos 3B e 7B em CPU & Sim & A definir (CPU/GPU no cluster) \\
Modelos 13B+ com GPU & Não viável no perfil simulado & Condicionado ao provisionamento ANTT \\
API REST FastAPI & Não & Sim \\
Deploy Kubernetes (Rancher) & Não (fora do escopo deste piloto) & Sim (ANTT + DeepFeed) \\
\bottomrule
\end{tabularx}
\end{table}

% =============================================================================
\section{Ambiente de desenvolvimento (simulação do deploy)}
% =============================================================================

A tabela abaixo descreve o hardware da estação usada no piloto. Os itens relevantes para a leitura deste documento são aqueles que \textbf{espelham o perfil de operação} (CPU-only, RAM finita, Ollama local). Detalhes de estação que não alteram esse perfil --- por exemplo, particularidades de instalação no WSL ou episódios pontuais de disco --- foram deliberadamente omitidos do diagnóstico de gargalos na Seção~\ref{sec:gargalos}.

\begin{table}[H]
\centering
\small
\caption{Perfil de hardware usado para simular o ambiente de operação}
\label{tab:hardware}
\begin{tabularx}{\textwidth}{l X}
\toprule
\textbf{Item} & \textbf{Observação} \\
\midrule
CPU & Intel i5-10300H --- suportou inferência quantizada 3B--7B, com latência elevada \\
RAM & 16~GiB totais; tipicamente 8--9~GiB disponíveis sob carga leve (ordem de grandeza compatível com o teto 7B em CPU) \\
Aceleração & Inferência \textbf{somente em CPU}, alinhada à restrição GETIC (sem GPU no ambiente permitido) e ao piloto CPU-only da arquitetura \\
Runtime LLM & Ollama (API compatível com OpenAI), modelo de referência \texttt{qwen2.5:7b} \\
\bottomrule
\end{tabularx}
\end{table}

\subsection{Correspondência com o ambiente de operação previsto}

\begin{table}[H]
\centering
\small
\caption{O que o desenvolvimento simulou versus o que fica para o cluster}
\label{tab:correspondencia}
\begin{tabularx}{\textwidth}{
  >{\raggedright\arraybackslash}p{4.2cm}
  >{\raggedright\arraybackslash}X
  >{\raggedright\arraybackslash}p{4.0cm}
}
\toprule
\textbf{Aspecto} & \textbf{Simulado no desenvolvimento} & \textbf{Deploy-alvo (arquitetura)} \\
\midrule
Ausência de GPU útil & Sim (CPU-only) & Sem GPU no piloto; GPU sob demanda futura \\
LLM local (Ollama) & Sim & Pod Ollama no Rancher \\
Faixa de modelo & 3B--7B quantizados & Piso técnico: 7~cores / 22~GB RAM, sem GPU \\
Dados no perímetro & Sim & Sim (sem saída obrigatória para API externa) \\
Orquestração K8s / Rancher & Não & Sim \\
Carga concorrente & Uma vaga local por processo (padrão); provedor externo sem fila neste backend & Mais vagas em CPU, ou GPU, dependem de cota da Agência \\
\bottomrule
\end{tabularx}
\end{table}

\subsection{Modelos compatíveis com o perfil simulado}

\begin{table}[H]
\centering
\small
\caption{Modelos locais avaliados no piloto}
\label{tab:modelos}
\begin{tabularx}{\textwidth}{
  >{\raggedright\arraybackslash}p{2.6cm}
  >{\centering\arraybackslash}p{2.0cm}
  >{\raggedright\arraybackslash}p{3.2cm}
  >{\raggedright\arraybackslash}X
}
\toprule
\textbf{Modelo} & \textbf{Disco} & \textbf{RAM aproximada} & \textbf{Papel no piloto} \\
\midrule
\texttt{llama3.2:3b} & $\sim$2,0~GB & Menor & Baseline rápido; capacidade limitada \\
\texttt{qwen2.5:7b} & $\sim$4,7~GB & $\sim$8--10~GiB & Maior modelo local praticável em CPU neste perfil \\
Modelos 13B+ & --- & Acima do teto do perfil & \textbf{Inviáveis} sem GPU e mais memória no ambiente de operação \\
\bottomrule
\end{tabularx}
\end{table}

Em suma: sob o perfil de operação simulado (CPU-only, memória finita), a infraestrutura \textbf{não comporta} LLMs locais da classe desejável para igualar, de forma estável, a qualidade observada em provedores de nuvem em tarefas normativas densas, sem degradar latência e estabilidade.

% =============================================================================
\section{Gargalos pertinentes ao perfil de operação}
\label{sec:gargalos}
% =============================================================================

Os itens abaixo são exclusivamente aqueles que se transferem do ambiente de desenvolvimento para o \textbf{ambiente de deploy/operação simulado} (CPU-only, Ollama local, memória limitada). Incidentes ou atritos de estação que não caracterizam esse perfil foram removidos deste diagnóstico.

\begin{table}[H]
\centering
\small
\caption{Gargalos de infraestrutura pertinentes ao perfil de operação}
\label{tab:gargalos_infra}
\begin{tabularx}{\textwidth}{
  >{\centering\arraybackslash}p{0.8cm}
  >{\raggedright\arraybackslash}p{3.8cm}
  >{\raggedright\arraybackslash}X
}
\toprule
\textbf{ID} & \textbf{Gargalo} & \textbf{Descrição / efeito} \\
\midrule
I1 & Inferência CPU-only & Sem GPU útil no perfil permitido (GETIC / piloto da arquitetura); toda a geração ocorre na CPU, com latência de dezenas de segundos a minutos \\
I2 & RAM insuficiente para modelos maiores & Com memória na ordem de 16~GiB na estação (e piso técnico de 22~GB no cluster-alvo sem GPU), o teto prático fica em modelos $\sim$7B quantizados; 13B+ induzem swap/thrashing ou inviabilidade operacional \\
I3 & Contenção de memória no fluxo RAG & Aplicação (Streamlit/API), embeddings locais e o modelo Ollama competem pela mesma quota de RAM; consultas longas pressionam o teto do perfil \\
I4 & Ausência de GPU no provisionamento ANTT & Enquanto o cluster não oferecer acelerador e quotas adequadas, o perfil de operação permanece CPU-only; o desenvolvimento reproduz essa restrição de propósito \\
\bottomrule
\end{tabularx}
\end{table}

\begin{figure}[H]
\centering
\footnotesize
\resizebox{\textwidth}{!}{%
\begin{tikzpicture}[
  node distance=6mm and 8mm,
  box/.style={draw, rounded corners, align=center, inner sep=4pt, minimum height=9mm},
  edge/.style={-Latex, thick}
]
\node[box, fill=gray!10, text width=3.2cm] (req) {Perfil de operação\\ANTT (CPU-only /\\Ollama local)};
\node[box, fill=red!10, text width=3.2cm, right=of req] (hw) {\textbf{I1--I4}\\CPU + RAM limitada\\sem GPU};
\node[box, fill=orange!10, text width=3.2cm, right=of hw] (mod) {Teto prático:\\modelos 3B--7B};
\node[box, fill=blue!5, text width=3.2cm, right=of mod] (gap) {Modelos mais\\robustos (13B+)\\fora de alcance};
\draw[edge] (req) -- (hw);
\draw[edge] (hw) -- (mod);
\draw[edge] (mod) -- (gap);
\end{tikzpicture}%
}
\caption{Cadeia causal: o perfil de operação simulado limita o teto do LLM local.}
\label{fig:cadeia_infra}
\end{figure}

\subsection{O que não entra neste diagnóstico}

Para manter o documento alinhado à premissa de simulação do deploy, \textbf{não} são tratados como gargalos de operação do RAG-ANTT:

\begin{itemize}
  \item Atritos exclusivos da estação de desenvolvimento (por exemplo, particularidades de WSL, instalação sem privilégios elevados ou episódios pontuais de espaço em disco);
  \item Otimizações internas do pipeline RAG (packing de contexto, prompts, retrieval) --- fora do escopo deste relatório de infraestrutura;
  \item Ausência, neste piloto, do deploy Kubernetes/Rancher propriamente dito (previsto na arquitetura, não requerido para validar o perfil CPU-only do LLM).
\end{itemize}

\subsection{Relação com o ambiente permitido pela Agência}

A decisão de operar LLM \textbf{local} atende a requisitos típicos de compliance (dados regulatórios no perímetro, previsibilidade de custo, independência de APIs externas). Contudo, a qualidade da geração local escala com o tamanho do modelo e com a disponibilidade de aceleração por GPU.

No perfil simulado --- CPU-only, memória finita, e produção tratada sob a mesma restrição até o provisionamento de GPU no Rancher --- a DeepFeed ficou limitada a validar modelos \textbf{3B e 7B}. Isso não é uma escolha arbitrária de produto: é a consequência direta de operar sob as mesmas restrições de infraestrutura esperadas no ambiente de deploy/operação atual da Agência.

Sem alteração desse teto (nós com GPU, quotas de RAM adequadas e runtime com aceleração no cluster), \textbf{não é tecnicamente responsável} comprometer o piloto ou a produção com modelos locais maiores, sob risco de instabilidade, latência inaceitável e contenção de memória.

\subsection{Atendimento simultâneo}

O backend distingue dois caminhos de geração. Eles não compartilham a mesma fila.

\textbf{Ollama em CPU.} A variável \texttt{RAG\_VAGAS\_GERACAO\_LOCAL} define quantas gerações locais este processo mantém ao mesmo tempo. O padrão do piloto é \textbf{1}: a segunda pessoa é aceita e espera a vez, no mesmo processo, até a primeira terminar. O código já admite um inteiro maior (teto interno 8). Subir esse número deixa o backend pronto para mais de uma inferência em CPU, desde que a Agência provisione núcleo e memória e o processo Ollama aceite mais de uma geração (\texttt{OLLAMA\_NUM\_PARALLEL} no serviço local, compatível com a mesma quantidade). Sem essa cota, vagas extras apenas disputam a mesma CPU e a mesma RAM do modelo 7B. Várias réplicas ou uma GPU continuam fora deste ciclo e dependem de capacidade no Rancher.

\textbf{Provedor externo (OpenAI ou DeepSeek).} Essas chamadas \textbf{não} entram na fila local. O processamento paralelo é por conta do serviço: várias pessoas podem gerar ao mesmo tempo até o limite de taxa e de crédito da conta. A busca nos documentos permanece no perímetro; o pedido de redação, com o trecho de contexto, sai para a internet. No perfil de produção fechado com a Agência esse caminho permanece desligado (\texttt{RAG\_LLM\_CLOUD\_FALLBACK=false}).

% =============================================================================
\section{Implicações práticas no perfil de operação}
% =============================================================================

\begin{itemize}
  \item Modelos \textbf{7B} em CPU são o teto prático do perfil simulado; já exigem cuidado com a quota de memória compartilhada pelo fluxo RAG;
  \item Modelos \textbf{13B ou superiores} são considerados inviáveis nesse perfil sem GPU e sem aumento de RAM no ambiente de operação;
  \item A primeira carga do modelo 7B no piloto levou da ordem de $\sim$1~minuto em pergunta curta; consultas RAG longas aumentam ainda mais o tempo de resposta --- latência coerente com inferência CPU-only;
  \item Enquanto a infraestrutura provisionada permanecer CPU-only, o LLM local deve ser dimensionado para o teto 3B--7B, e expectativas de qualidade devem ser calibradas em relação a esse teto --- ou a Agência deve autorizar/provisionar GPU e memória compatíveis com modelos mais robustos;
  \item Várias pessoas podem consultar ao mesmo tempo. No Ollama local o padrão é uma geração por vez neste processo (\texttt{RAG\_VAGAS\_GERACAO\_LOCAL}); em provedor externo o paralelismo é do serviço, e não desta fila.
\end{itemize}

% =============================================================================
\section{Como reproduzir o ambiente do piloto}
% =============================================================================

\begin{lstlisting}
# 1) Servidor Ollama (CPU) --- perfil de operacao simulado
export PATH="$HOME/.local/bin:$PATH"
export LD_LIBRARY_PATH="$HOME/.local/ollama/lib:$LD_LIBRARY_PATH"
ollama serve
ollama list   # llama3.2:3b e qwen2.5:7b

# 2) Aplicacao
cd RAG-ANTT
# Fila local do Ollama. Padrao 1. OpenAI/DeepSeek ignoram esta variavel.
# export RAG_VAGAS_GERACAO_LOCAL=1
streamlit run antt_rag_unified.py

# 3) UI
# Sidebar -> Servico: Local (Ollama) - CPU
# Modelo: qwen2.5:7b
\end{lstlisting}

% =============================================================================
\section{Conclusão}
% =============================================================================

O piloto demonstrou que é \textbf{viável} operar LLM local via Ollama no fluxo Streamlit do RAG-ANTT sob o \textbf{perfil de deploy/operação simulado} (CPU-only, sem GPU). Os gargalos determinantes para não adotar modelos locais mais robustos são os que se transferem para esse perfil:

\begin{enumerate}[label=\textbf{\arabic*.}]
  \item Inferência \textbf{CPU-only}, alinhada à restrição do ambiente permitido pela ANTT até o provisionamento de GPU;
  \item \textbf{RAM} insuficiente, nesse perfil, para carregar com estabilidade modelos acima da faixa $\sim$7B quantizados;
  \item Contenção de memória entre aplicação, embeddings e o próprio modelo no fluxo RAG;
  \item Persistência da ausência de GPU no desenho atual de provisionamento até nova decisão da Agência.
\end{enumerate}

Portanto, a limitação ao uso de LLMs locais mais robustos decorre das \textbf{condições de infraestrutura do ambiente de operação que o desenvolvimento visou simular}, e não de uma recusa a priori de modelos maiores nem de atritos locais de estação. A evolução da qualidade do LLM local na ANTT depende, em primeiro lugar, do \textbf{provisionamento} (GPU e/ou mais RAM no cluster Rancher) alinhado aos requisitos de compliance e desempenho da Agência.

O mesmo vale para o atendimento simultâneo. O backend já separa os caminhos: no Ollama, \texttt{RAG\_VAGAS\_GERACAO\_LOCAL} (padrão uma vaga) segura a fila local e pode ser elevado quando houver CPU e RAM; em OpenAI ou DeepSeek não há essa fila, e o processamento paralelo fica por conta do serviço. O perfil de produção da Agência segue no modelo local.

% =============================================================================
\vspace{1cm}
\begin{center}
\rule{0.5\textwidth}{0.4pt}
\end{center}
\vspace{0.5cm}

\noindent\textit{Documento vivo: atualizar quando mudarem o perfil de infraestrutura ANTT, o hardware usado na simulação do deploy ou o teto de modelo local suportado. Última revisão: setembro de 2026 (premissa: desenvolvimento como simulação do perfil de operação; gargalos restritos a esse perfil; fila local de geração e paralelismo do provedor externo).}

\end{document}
